\section{Discussão e limitações}
A contribuição ilustrada pelo quadro é separar enunciado original, regra executável e limite de verificação; essa separação torna a análise revisável sem atribuir decisões individuais na ausência de registros de autoria. Não se calculam percentuais de fidelidade com unidades escolhidas para demonstração, pois o corpus completo ainda não foi classificado.

As margens de 3 cm na parte superior e esquerda, 2 cm na inferior e direita, fonte de 12 pontos e entrelinhas de 1,5 constituem uma adaptação gráfica de trabalho acadêmico; a norma específica de artigo e as futuras exigências editoriais devem ser conferidas separadamente \citep{abnt2024,ufes2026}. Resumos, quadros e referências usam espaço simples para preservar sua organização e legibilidade \citep{abnt2024}.

O template canônico consultado foi publicado em 2018, anterior às atualizações de resumos, citações e referências identificadas em 2021, 2023 e 2025; o modelo presente aplica os requisitos lidos e registra a necessidade de revisão integral das normas ausentes no corpus \citep{abntex2018,ufes2026}. As normas antigas de sumário, índice e lombada presentes no acervo não são transferidas automaticamente para um artigo em PDF.

\section{Considerações finais}
A demonstração oferece uma estrutura modular com metadados, resumo, introdução, fundamentação, método, resultados, discussão e referências; sua utilidade é editorial e não equivale à validação do estudo ou à aceitação por periódico \citep{abntex2018}. Para reutilizá-la, o pesquisador deve substituir o corpus, revisar as fontes e conferir a correspondência entre cada conclusão e sua evidência.

\section*{\centering Transparência}
A organização e a redação deste exemplo receberam assistência de ChatGPT/Codex, conforme a proveniência registrada no projeto; autoria humana, revisão intelectual e atribuições individuais devem ser conferidas pelos pesquisadores. Os dados técnicos são extraídos do registro piloto e não representam dados simulados de aprendizagem.

\section*{\centering Rastreabilidade das referências técnicas}
O modelo canônico consultado existe em código-fonte público e em PDF de seis páginas, com identificação da versão 1.9.7 de 2018; sua leitura permite conferir os comandos e as regras exemplificadas, mas não comprova atualização a todas as normas posteriores \citep{abntex2018}. O código foi conferido no GitHub, com blob \nolinkurl{9c38dc56d0d7b4fd62267ba2394afa484049ce52}. O PDF externo está disponível em \url{https://ctan.math.washington.edu/tex-archive/macros/latex/contrib/abntex2/doc/examples/abntex2-modelo-artigo.pdf} \citep{abntex2018}.

A página de normalização da UFES foi lida e informa atualização em 10 de julho de 2026, distinguindo as edições atuais das normas dos manuais antigos; essa página é uma fonte institucional de identificação das edições, e não substitui a leitura integral dos textos normativos \citep{ufes2026}. Os PDFs NBR 6028:2021, 10520:2023 e 14724:2024 utilizados nesta demonstração foram lidos no acervo privado fornecido pelo pesquisador; eles são normas técnicas, e não artigos científicos. A NBR 6023:2025 e a NBR 6022:2018 foram identificadas em fonte institucional, sem leitura integral nesta etapa, portanto não se afirma conformidade integral com elas \citep{ufes2026}.
